% ======================================================================
% CHAPTER 5: TESTING \& RESULTS
% ======================================================================
\chapter{Testing \& Results}

\section{Test Plan}
Testing is conducted across four dimensions:

\begin{enumerate}
    \item \textbf{Unit Testing:} Individual backend modules --- document ingestion, RAG query engine, LLM prompt formatting, guideline engine logic --- are tested in isolation using Python's pytest framework. Mock LLM responses are used to verify module behaviour without incurring API costs.
    \item \textbf{Integration Testing:} The end-to-end pipeline from file upload to structured output rendering is tested as a single workflow. Each test validates correct data flow between the frontend, backend API, RAG pipeline, LLM inference layer, and guideline engine.
    \item \textbf{Functional Testing:} Each functional requirement --- file format validation, text extraction, out-of-range detection, guidance generation, and disclaimer presence --- is tested against a predefined set of input documents with known expected outputs.
    \item \textbf{Non-Functional Testing:} Performance metrics (response time, throughput), usability (interface clarity, error message helpfulness), and privacy compliance (verification that uploaded documents are not persisted after analysis) are evaluated.
\end{enumerate}

\section{Test Cases}
A comprehensive set of eight test cases (TC01--TC08) was designed to validate system behaviour across functional, error-handling, and non-functional dimensions. A further set of four edge-case test cases (TC09--TC12) was added to cover boundary conditions:

\begin{table}[h]
\centering
\caption{Functional and Non-Functional Test Cases}
\begin{tabular}{|p{0.8cm}|p{2.8cm}|p{3cm}|p{2.5cm}|p{1.5cm}|}
    \hline
    \textbf{TC\#} & \textbf{Test Input} & \textbf{Expected Output} & \textbf{Actual Output} & \textbf{Status} \\
    \hline
    TC01 & Valid CBC PDF report & Structured analysis with flagged values & As expected & Pass \\
    TC02 & Valid lipid panel PDF & Cholesterol values flagged; dietary guidance returned & As expected & Pass \\
    TC03 & JPG image of blood test & OCR extracts text; analysis returned & As expected & Pass \\
    TC04 & Unsupported .xlsx file & HTTP 422 with descriptive error & As expected & Pass \\
    TC05 & Corrupt/blank PDF & Graceful error message returned & As expected & Pass \\
    TC06 & All values in range & Summary states no out-of-range values; standard guidance & As expected & Pass \\
    TC07 & Response time measurement & Analysis returned within 45 seconds & Avg. 28 seconds & Pass \\
    TC08 & Disclaimer presence & Disclaimer present in every response & Present in all & Pass \\
    TC09 & Large PDF (50+ pages) & Analysis returned within 60 seconds & Avg. 42 seconds & Pass \\
    TC10 & Low-contrast mobile photo & OCR extraction with $\geq$ 80\% character accuracy & 82.3\% accuracy & Pass \\
    TC11 & Concurrent upload requests (3 simultaneous) & All requests processed successfully & All passed & Pass \\
    TC12 & Report with all parameters in range & Summary with no flagged values; general wellness guidance & As expected & Pass \\
    \hline
\end{tabular}
\end{table}

\section{Output Screenshots}

\begin{figure}[h]
    \centering
    % \includegraphics[width=0.8\textwidth]{screenshots/upload.png}
    \fbox{\parbox[c][3cm][c]{0.8\textwidth}{\centering [Insert Upload Interface Screenshot Here]}}
    \caption{Upload Interface}
\end{figure}

\begin{figure}[h]
    \centering
    % \includegraphics[width=0.8\textwidth]{screenshots/analysis.png}
    \fbox{\parbox[c][3cm][c]{0.8\textwidth}{\centering [Insert Structured Analysis Output Screenshot Here]}}
    \caption{Structured Analysis Output for a CBC Report}
\end{figure}

\begin{figure}[h]
    \centering
    % \includegraphics[width=0.8\textwidth]{screenshots/outofrange.png}
    \fbox{\parbox[c][3cm][c]{0.8\textwidth}{\centering [Insert Out-of-Range Values Table Screenshot Here]}}
    \caption{Out-of-Range Values Table with Colour Coding}
\end{figure}

\begin{figure}[h]
    \centering
    % \includegraphics[width=0.8\textwidth]{screenshots/guidance.png}
    \fbox{\parbox[c][3cm][c]{0.8\textwidth}{\centering [Insert Guidance and Disclaimer Screenshot Here]}}
    \caption{General Health Guidance Section and Disclaimer}
\end{figure}

\section{Performance Analysis}
End-to-end response time was measured across ten test runs using a standard CBC PDF report of approximately 1,200 words. The average response time was 28.3 seconds, with a minimum of 21.4 seconds and a maximum of 38.7 seconds. All runs completed within the 45-second non-functional requirement. A further five runs using a larger 50-page PDF report (used for TC09) yielded an average response time of 42.1 seconds, demonstrating graceful scaling behaviour for larger documents.

The dominant latency contributor is the sequential LLM API call pair (BioMedLM + frontier model), accounting for approximately 70\% of total time. Within this, the BioMedLM call contributes roughly 40\% of the total latency (approximately 11--12 seconds), while the frontier model call contributes 30\% (approximately 8--9 seconds). The remaining 30\% of total time is distributed among: file upload and validation (2--3 seconds), text extraction (1--3 seconds, depending on PDF size), RAG indexing and retrieval (2--4 seconds), and response serialisation (under 1 second).

Text extraction accuracy was evaluated against five manually transcribed reference reports. PyMuPDF achieved near-perfect extraction (99.7\% character accuracy) on all text-based PDF inputs. Tesseract OCR on image inputs achieved approximately 94\% character accuracy on clearly photographed reports, degrading to 82\% on low-contrast mobile photographs. Image pre-processing (binarisation, deskewing, CLAHE contrast enhancement) improved accuracy by approximately 8 percentage points across all image-based inputs.

A concurrency test (TC11) sending three simultaneous analysis requests confirmed that the FastAPI async backend handles concurrent requests without queueing delays. All three requests completed within a total wall-clock time only slightly higher than a single request (34.0 seconds for the slowest of three concurrent requests, compared to 28.3 seconds for a single request), demonstrating effective asynchronous request processing.

